\chapterstar{Objetivo General}
Diseñar e implementar un sistema de gestión de inventario que, con una base de datos en PostgreSQL, una API desarrollada en FastAPI y una aplicación web desarrollada en React, realice el registro de control de existencias y ventas utilizando un lector de códigos de barras, y realice el monitoreo de las variaciones de temperatura y humedad con un microcontrolador ESP32 con sensor AHT10, para la Farmacia Selene.

\chapterstar{Objetivos Específicos}

\begin{enumerate}
\item Desarrollar una base de datos en PostgreSQL, para almacenar de forma estructurada la información de productos, niveles de inventario, usuarios, registros de ventas y mediciones obtenidas de los sensores para Farmacia Selene.

\item Desarrollar una API en Python utilizando FastAPI y SQLAlchemy para gestionar el registro de inventario, salidas de inventario, autenticación de usuarios y consulta de registros del sensor para publicarla en la plataforma Railway para su integración con la aplicación desarrollada en React.

\item Desarrollar el Frontend mediante una aplicación en JavaScript utilizando React para la gestión de inventario, mostrando tableros gráficos, reportes y alertas generadas a partir de los datos obtenidos del sensor para publicarla en la plataforma Vercel para su acceso público.

\item Desarrollar un programa en C++ para que el microcontrolador ESP32 adquiera los datos de temperatura y humedad mediante el sensor AHT10 y envíen esta información a la API a través de solicitudes HTTP.  

\item Diseñar e implementar un gabinete de plástico hecho por una impresora 3D con medidas de 12x5x5cm, bajo el software Bambu Studio para alojar el microcontrolador ESP32, el sensor AHT10, la fuente de alimentación y tres pantallas OLED SSD1306 de 128x64 píxeles, cada pantalla sera para una instalación segura en el área de medición de temperatura, humedad y estado de conexión de WiFi, integrando el lector de códigos de barras en la aplicación de React para optimizar el registro y actualización de existencias en Farmacia Selene.

\end{enumerate}

\thispagestyle{empty}
